\begin{abstract}
The Robertson problem is a three-species stiff chemical-kinetics initial-value problem. We compare explicit Euler (EE), classical fourth-order Runge--Kutta (RK4), and implicit Euler (IE) solved by Newton iteration. Along a tight reference trajectory, the two dynamically relevant Jacobian modes develop widely separated decay rates, giving a stiffness ratio of $1.5840\times10^5$ at $t=40$ and a fast-mode magnitude of $3392.79$. Frozen-Jacobian analysis predicts $h_{EE}\lesssim5.895\times10^{-4}$ and $h_{RK4}\lesssim8.209\times10^{-4}$; these are local diagnostics rather than nonlinear stability theorems. The observed convergence orders are $1.079$, $4.268$, and $0.99996$ for EE, RK4, and IE, respectively, while the largest stable tested steps for EE and RK4 are $5.4159\times10^{-4}$ and $8.3088\times10^{-4}$. The reference endpoint is $y_2(40)=9.185535\times10^{-6}$, supported to ten significant digits by tolerance repetition and two independent solvers. A stability-aware adaptive RK4 variant caps the error-based step proposal using the local reduced-Jacobian stability scale. At the three looser tolerances this removes rejected trials and reduces RHS work by approximately $22$--$24\%$, whereas at the strictest tolerance the cap only adds spectral-evaluation overhead. We also compare work at matched error, test conservation and non-negativity as diagnostics, examine the effects of Newton and adaptive tolerances, and benchmark six \texttt{solve\_ivp} solvers at a common tolerance: the non-stiff solvers need $1.4\times10^{5}$--$2.4\times10^{5}$ right-hand-side evaluations where the stiff solvers need $405$--$752$, a two-order-of-magnitude gap at matched accuracy.
\end{abstract}
